\section{Conclusions}

In this work, I analyzed the RFI environment at the DRAO using a one-hour segment of data from the DRAO RFI monitor and estimated the number of detectable H II lines. As part of the analysis I found that each calibration segment necessitated an independent RFI label to account for changes in the gain factor. I applied spectral kurtosis to integrated power and, after proper calibration, found that it closely matched theory. The generated RFI labels were compared against known RFI events and found strong agreement. We considered a frequency bin clean provided that it was RFI-free in at least 50\% of the time segments, and we applied a Gaussian-weighted average over frequency bins to classify whether a given RRL would be visible. I found that a total of 79 RRLs in the range of 300 -- 1500 MHz are detectable and 37 are contaminated by RFI. The 79 detectable lines are used as the basis for all future detection analysis.\\

I derived the theoretical point source sensitivity and brightness temperature sensitivity of CHORD and found that intrinsic sensitivity increases at high declination and at low frequency. Radio background adds linearly to the system temperature and is strongest at low frequency and along the Galactic plane. I found that sensitivity degrades at low frequency due to the strong background and that the overall sensitivity maps are highly spatially dependent and frequency dependent. Rather then estimating a single CHORD sensitivity value, this work presents sky maps which provide sensitivity for any given frequency or sky location.\\

I showed how an inverse-variance-weighted average can be used to maximize the signal-to-noise ratio of multiple RRLs having different sensitivities. I explored how the optimal RRL stacking order varies by sky location due to the changing SNR per line. I analyzed both velocity and angular smoothing to show that lines can be stacked more effectively at the cost to angular precision. As with sensitivity, the optimal stacking strategy depends on sky location and on scientific goals, namely which of sensitivity, velocity resolution, or angular resolution is most valuable.\\

I modelled adding dishes to the main CHORD array to create the Extension CHORD array. We identified 8 locations as candidate sites and modelled all combinations of 1 or 2 locations, for 40 total combinations. I found that the Synthesis Telescope track in combination with the disturbed land south of CHIME represented the optimal locations for constructing Extension CHORD. While these two locations do not represent the longest possible baselines, they work together to produce excellent \UV\ coverage and a symmetric synthesized beam without any distortions. I found that adding 32 dishes to the recommended locations could improve angular resolution by 44\% while reducing sensitivity by 25\%, or could improve angular resolution by 63\% for a loss in sensitivity of 200\%.\\

Finally, I estimated CHORD's detection ability of real H II regions by simulating the CHORD response on data published by the SIGGMA survey. In the worst case, CHORD could match SIGGMA SNR levels with a common 6 arcminute angular resolution after 4 observing days, and could improve angular resolution to 5 arcminutes with matching SNR levels after 9 days. In other cases CHORD can detect the bright H II regions measured by SIGGMA after just one day and in most cases using only one RRL. In conclusion, I've shown that an RRL survey with CHORD could both improve detections of known H II regions and likely measure new ones, building a more complete picture of the ionized interstellar medium.\\

